% ---------------------------------------------------------------------------
% Abstract and keywords
% ---------------------------------------------------------------------------

\chapter*{Abstract}

ISOGEO's Scan product explores various data sources in order to identify the
resources they contain and extract the information required for their
documentation. Its operation relies on several processes capable of analysing
databases, geodatabases and geospatial files from different technical
environments.

Some of these processes were historically implemented using FME workbenches.
Although FME provides numerous data integration and transformation features,
its use introduced an external dependency into the Scan execution chain. The
main objective of this apprenticeship was therefore to progressively replace
some of these workflows with a Python solution integrated into the product and
capable of running independently.

The first stage of the work consisted in analysing the existing FME workflows
in order to identify their inputs, outputs, transformations and business rules.
Python components were then developed to connect to several data sources,
including PostgreSQL/PostGIS, Oracle, SQL Server and Esri enterprise
geodatabases. These developments extract information about tables, fields,
coordinate reference systems, spatial extents, geometries and various technical
metadata required by Scan.

The GDAL/OGR library was used in particular to analyse vector and raster
geospatial data sources. Specific implementations were also introduced to
handle the characteristics of each database management system, invalid data,
connection errors and behavioural differences between formats.

Particular attention was paid to the industrialisation of the solution. The
developed components were integrated into a modular architecture, supported by
structured error management and prepared for the generation of a standalone
Windows executable using PyInstaller. A regression testing approach was also
designed to automatically compare the results produced by the former FME
workflows with those obtained from the new Python implementation.

The work carried out contributes to reducing Scan's dependency on FME,
improving the maintainability of its processing components and facilitating
their integration into ISOGEO's software environment. This apprenticeship also
enabled me to strengthen my skills in Python development, geospatial data
processing, spatial databases, software packaging and technical solution
validation.

\vspace{0.7cm}

\noindent
\textbf{Keywords:}
geomatics; Scan; ISOGEO; Python; FME; GDAL/OGR; metadata; spatial databases;
PyInstaller; regression testing.
